% ============================================================================
%  SECCIÓN 3.20: COMPONENTES REUTILIZABLES
% ============================================================================

\section{Componentes reutilizables}

La reutilizaci\'on de componentes es un principio fundamental del desarrollo
de software que reduce la duplicaci\'on de c\'odigo, facilita el mantenimiento
y garantiza la consistencia visual en toda la aplicaci\'on
\citep{sommerville2011}. En el presente proyecto se implementaron los
siguientes componentes y patrones reutilizables.

\subsection{Componente AuthGuard (vigilante de autenticaci\'on)}

El componente \texttt{AuthGuard} es un componente de alto orden que envuelve
toda la aplicaci\'on y controla el acceso a las pantallas seg\'un el estado de
autenticaci\'on del usuario. Implementa la l\'ogica de redireccionamiento
autom\'atico: si el usuario no est\'a autenticado y no se encuentra en una
pantalla de autenticaci\'on, es redirigido a la pantalla de splash; si est\'a
autenticado pero se encuentra en una pantalla de autenticaci\'on, es redirigido
a la pantalla principal.

\begin{lstlisting}
export function AuthGuard({ children }: { children: React.ReactNode }) {
  const { session, loading } = useAuth();
  const segments = useSegments();
  const router = useRouter();

  useEffect(() => {
    if (loading) return;

    const inAuthGroup = segments[0] === 'auth';
    const inSplash = segments[0] === 'splash';

    if (!session && !inAuthGroup && !inSplash) {
      router.replace('/splash');
    } else if (session && (inAuthGroup || inSplash)) {
      router.replace('/(tabs)');
    }
  }, [session, loading, segments, router]);

  if (loading) {
    return (
      <View style={styles.loadingContainer}>
        <ActivityIndicator size="large" color={colors.primary} />
      </View>
    );
  }

  return <>{children}</>;
}
\end{lstlisting}

Este componente se integra en el layout ra\'iz de la aplicaci\'on
(\texttt{\_layout.tsx}), envolviendo a todos los componentes hijos. La
ventaja de este patr\'on es que centraliza la l\'ogica de protecci\'on de rutas
en un solo lugar, evitando duplicar c\'odigo en cada pantalla.

\begin{figure}[H]
  \centering
  \caption{Diagrama de flujo del componente AuthGuard}
  \label{fig:authguard-flow}
  \contenidoConFuente{%
  \begin{tikzpicture}[
      nodo/.style={rectangle, draw=black!70, rounded corners=2pt,
                   minimum width=3.0cm, minimum height=0.9cm,
                   align=center, fill=gray!10, font=\small},
      decision/.style={diamond, draw=black!70, rounded corners=2pt,
                       minimum width=2.8cm, minimum height=1.0cm,
                       align=center, fill=orange!10, font=\small},
      flecha/.style={-{Stealth[length=2.5mm]}, thick}
    ]
    \node[decision] (check) at (0,0) {\'Es autenticado?};
    \node[nodo] (splash) at (-3.5,-2) {Redirigir a Splash};
    \node[nodo] (tabs) at (3.5,-2) {Redirigir a Tabs};
    \node[nodo] (auth) at (3.5,2) {Permitir acceso};

    \draw[flecha] (check) -- node[below left] {No} (splash);
    \draw[flecha] (check) -- node[below right] {S\'i} (auth);
    \draw[flecha] (tabs) -- node[above right] {S\'i} (check);
  \end{tikzpicture}%
  }{Elaboraci\'on propia en base al c\'odigo fuente del componente.}
\end{figure}

\subsection{Contexto de autenticaci\'on (AuthContext)}

El \texttt{AuthContext} es un contexto de React que centraliza el estado de
autenticaci\'on de la aplicaci\'on, proporcionando a cualquier componente hijo
acceso al usuario actual, la sesi\'on y las funciones de inicio de sesi\'on,
registro y cierre de sesi\'on. El contexto gestiona la suscripci\'on a cambios
de autenticaci\'on en tiempo real a trav\'es de la API de Supabase.

\begin{lstlisting}
interface AuthContextType {
  session: Session | null;
  user: User | null;
  loading: boolean;
  signIn: (email: string, password: string)
    => Promise<{ error?: string }>;
  signUp: (email: string, password: string, fullName: string)
    => Promise<{ error?: string }>;
  signOut: () => Promise<void>;
}
\end{lstlisting}

El contexto se define una sola vez en el layout ra\'iz de la aplicaci\'on y
est\'a disponible para todas las pantallas y componentes a trav\'es del hook
\texttt{useAuth()}.

\subsection{Sistema de tema visual}

El sistema de tema visual se define en el directorio \texttt{src/theme/} y
centraliza los colores, espaciados, tama\~nos de fuente y radios de borde de
la aplicaci\'on. Este patr\'on garantiza que todas las pantallas mantengan una
identidad visual coherente.

\subsubsection{Paleta de colores}

La paleta de colores se define en el archivo \texttt{colors.ts} y utiliza un
sistema de tokens de diseño que asocia cada color a un nombre sem\'antico:

\begin{lstlisting}
export const colors = {
  primary: '#2563EB',
  primaryDark: '#1D4ED8',
  primaryLight: '#93C5FD',

  secondary: '#10B981',
  secondaryDark: '#059669',
  secondaryLight: '#6EE7B7',

  background: '#F8FAFC',
  surface: '#FFFFFF',
  surfaceVariant: '#F1F5F9',

  text: '#0F172A',
  textSecondary: '#64748B',
  textInverse: '#FFFFFF',

  border: '#E2E8F0',
  borderLight: '#F1F5F9',

  error: '#EF4444',
  errorLight: '#FCA5A5',
  success: '#22C55E',
  warning: '#F59E0B',
} as const;
\end{lstlisting}

La paleta se inspira en la escala de colores de Tailwind CSS, lo que permite
mantener un estilo visual moderno y consistente. Los colores se agrupan en
categor\'ias: primarios (azul), secundarios (verde), superficies (grises),
texto, bordes y estados (error, \'exito, advertencia).

\subsubsection{Sistema de espaciado}

El sistema de espaciado se define en el archivo \texttt{spacing.ts} y
establece una escala de valores fijos que se aplican de forma consistente en
todos los componentes:

\begin{lstlisting}
export const spacing = {
  xs: 4, sm: 8, md: 16, lg: 24, xl: 32, xxl: 48,
} as const;

export const borderRadius = {
  sm: 4, md: 8, lg: 12, xl: 16, full: 9999,
} as const;

export const fontSize = {
  xs: 12, sm: 14, md: 16, lg: 18, xl: 20, xxl: 24, title: 32,
} as const;
\end{lstlisting}

Este sistema de tokens se exporta desde el archivo \texttt{index.ts} del
directorio \texttt{theme/} y se importa en cada pantalla mediante una sola
importaci\'on, simplificando el acceso a las constantes de dise\~no.

\subsection{Tipos de datos compartidos}

Los tipos de datos se definen en el directorio \texttt{src/types/} e incluyen
las interfaces para los modelos de dominio de la aplicaci\'on: usuarios,
escenarios, deportes, eventos, favoritos e im\'agenes de escenarios. Estos
tipos se utilizan en toda la aplicaci\'on para garantizar la consistencia de
los datos.

\begin{lstlisting}
export interface User {
  id: string;
  email: string;
  full_name: string;
  role: 'admin' | 'gestor' | 'asistente';
  created_at: string;
  updated_at: string;
}

export interface Scenario {
  id: string;
  nombre: string;
  tipo: string;
  descripcion: string;
  capacidad: number;
  direccion: string;
  latitud: number;
  longitud: number;
  estado: string;
  created_by: string | null;
  created_at: string;
  updated_at: string;
}
\end{lstlisting}
